% Part: lambda-calculus
% Chapter: introduction
% Section: reduction

\documentclass[../../../include/open-logic-section]{subfiles}

\begin{document}

\olfileid{lam}{int}{red}
\olsection{λ-项的归约}

我们能用λ-项做什么呢？简化它们。如果 $M$ 和 $N$ 是任意λ-项而 $x$ 是任意变元，
我们可以用 $\Subst{M}{N}{x}$ 表示将 $N$ 代入 $M$ 中的~$x$ 后得到的结果，
在此过程中，需先重命名 $M$ 中那些在代入后可能与~$N$ 的自由变元发生冲突的约束变元。
例如，
\[
\Subst{(\lambd[w][xxw])}{yyz}{x} = \lambd[w][(yyz)(yyz)w].
\]

\begin{digress}
替换的其他记法有 $[N/x]M$、$[x/N]M$ 以及 $M[x/N]$。当心！
\end{digress}

直观上说，$(\lambd[x][M])N$ 和 $\Subst{M}{N}{x}$ 意义相同；
将前一项替换为后一项的操作称为\emph{$\beta$-收缩}。
$(\lambd[x][M])N$ 称为\emph{可约式}，而 $\Subst{M}{N}{x}$ 称为其\emph{合约项}。
一般地，如果通过对项 $P$ 的某个子项进行 $\beta$-收缩，可以将 $P$ 变为~$P'$，
我们就说 $P$ \emph{在一步内 $\beta$-归约}到 $P'$，并记作 $P \redone P'$。
如果从 $P$ 出发，经过若干次（可能零次）一步归约能得到~$P'$，
那么 $P$ \emph{$\beta$-归约}到~$P'$；记作 $P \red P'$。
一个不能再进行任何 $\beta$-归约的项称为\emph{$\beta$-不可约的}，或\emph{$\beta$-正规的}。
当上下文清楚时，我们就直接说“归约”而不说“$\beta$-归约”等。

让我们来看几个例子。
\begin{enumerate}
\item 我们有
\begin{align*}
(\lambd[x][xxy]) \lambd[z][z] & \redone (\lambd[z][z])(\lambd[z][z]) y \\
& \redone (\lambd[z][z]) y \\
& \redone y.
\end{align*}
\item “简化”一个项可能使其变得更复杂：
\begin{align*}
(\lambd[x][xxy])(\lambd[x][xxy]) & \redone (\lambd[x][xxy])(\lambd[x][xxy])y \\
& \redone (\lambd[x][xxy])(\lambd[x][xxy])yy \\
& \redone \dots
\end{align*}
\item 它也可能让项保持不变：
\[
(\lambd[x][xx])(\lambd[x][xx]) \redone (\lambd[x][xx])(\lambd[x][xx]).
\]
\item 此外，有些项可以按不止一种方式归约；例如，
\[
(\lambd[x][(\lambd[y][yx]) z)] v \redone (\lambd[y][yv]) z
\]
这是通过收缩最外层的应用得到的；而
\[
(\lambd[x][(\lambd[y][yx]) z)] v \redone (\lambd[x][zx]) v
\]
则是通过收缩最内层的应用得到的。不过请注意，在这个例子中，两项进一步归约到了同一个项 $zv$。
\end{enumerate}

上一个例子中的最终结果并非巧合，而是体现了λ-演算一个深刻而重要的性质，即所谓的“Church--Rosser性质”。

\end{document}